\documentclass[12pt,a4paper]{article}
\usepackage[utf8]{vietnam}
\usepackage[top=2cm, bottom=2cm, left=1.5cm, right=1cm]{geometry}
\usepackage{amsmath,amsfonts,amssymb}
\usepackage{tabularx,longtable,multirow}
\usepackage{array}
\usepackage{fancyhdr}
\usepackage{enumitem}
\usepackage{xcolor}

% Cấu hình cột bảng để không bị tràn lề
\newcolumntype{L}[1]{>{\raggedright\arraybackslash}p{#1}}
\newcolumntype{C}[1]{>{\centering\arraybackslash}p{#1}}

% Cấu hình Header/Footer
\pagestyle{fancy}
\fancyhf{}
\rhead{\small GV: Nguyễn Văn Sang}
\lhead{\small Kế hoạch giáo dục 2025-2026 - Môn Toán 11}
\cfoot{\thepage}

\begin{document}

% --- PHẦN TIÊU NGỮ (TRANG 1) ---
\noindent
\begin{minipage}[t]{0.5\textwidth}
    \textbf{TRƯỜNG: THPT NGUYỄN HỮU CẢNH}\\
    \textbf{TỔ: Toán}\\
    \textbf{Họ và tên giáo viên: Nguyễn Văn Sang}
\end{minipage}
\begin{minipage}[t]{0.5\textwidth}
    \centering
    \textbf{CỘNG HÒA XÃ HỘI CHỦ NGHĨA VIỆT NAM}\\
    \textbf{Độc lập - Tự do - Hạnh phúc}\\
    ----------------------------
\end{minipage}

\vspace{0.8cm}
\begin{center}
    \Large \textbf{KẾ HOẠCH GIÁO DỤC CỦA GIÁO VIÊN} \\
    \large \textbf{(Năm học 2025 - 2026)}
\end{center}

\section*{I. THÔNG TIN CÁ NHÂN}
\begin{itemize}[noitemsep]
    \item Họ và tên: Nguyễn Văn Sang \quad | \quad Năm sinh: 1987
    \item Năm vào ngành: 2010 \quad | \quad Chức vụ: Giáo viên
    \item Giảng dạy: Toán \quad | \quad Công tác kiêm nhiệm: Không
\end{itemize}

\section*{II. KẾ HOẠCH CÁ NHÂN}
\subsection*{1. Những căn cứ xây dựng kế hoạch}
\begin{itemize}[label=-, noitemsep]
    \item Công văn 4612/BGDĐT-GDTrH ngày 03/10/2017 về hướng dẫn chương trình GDPT.
    \item Quyết định số 3089/QĐ-UBND ngày 08/8/2025 của UBND TP. Hồ Chí Minh.
    \item Văn bản số 4977/SGDĐT-GDTrH ngày 13/8/2025 của Sở Giáo dục và Đào tạo TP. HCM.
    \item Kế hoạch số 4818/KH-SGDĐT ngày 07/8/2025 về bồi dưỡng thường xuyên.
    \item Căn cứ vào năng lực cá nhân và yêu cầu nhiệm vụ được giao.
\end{itemize}

\subsection*{2. Mục tiêu chung}
- Hoàn thành tốt nhiệm vụ được giao. Đáp ứng Chuẩn nghề nghiệp, phát triển bản thân.

\subsection*{3. Nội dung}
\textbf{3.1. Đặc điểm tình hình:}\\
- \textit{Thuận lợi:} Có tinh thần trách nhiệm cao, nhiệt tình. Trình độ chuyên môn vững vàng, ứng dụng tốt CNTT.\\
- \textit{Khó khăn:} CSVC còn thiếu thốn. Cả 3 khối 10, 11, 12 đều học chương trình mới nên nhiều thứ cần chuẩn bị.

\textbf{3.1.2 Công việc được giao:}\\
Giảng dạy môn Toán các lớp: 11B5, 11B15.

\textbf{3.2.2 Giáo dục đạo đức, tư tưởng:} Giáo dục HS lòng yêu nước, tôn trọng pháp luật, nhân ái và tự trọng.

\newpage
% --- III. PHÂN PHỐI CHƯƠNG TRÌNH CHI TIẾT ---
\section*{III. KẾ HOẠCH DẠY HỌC (PHÂN PHỐI CHƯƠNG TRÌNH)}
\textbf{HỌC KỲ I}
\begin{longtable}{|C{1.1cm}|C{0.5cm}|L{4.2cm}|C{0.5cm}|L{4.2cm}|C{0.5cm}|L{2.5cm}|}
\hline
\textbf{Tuần} & \textbf{T} & \textbf{ĐS-GT-TK} & \textbf{T} & \textbf{HHKG} & \textbf{T} & \textbf{Chuyên đề} \\ \hline
\endhead

% TUẦN 1
\multirow{2}{*}{\begin{tabular}[c]{@{}c@{}}1\\(8/9-\\14/9)\end{tabular}} & 1 & \textbf{Góc lượng giác. Giá trị lượng giác.} \newline \textit{Yêu cầu:} Nhận biết khái niệm góc LG, số đo, hệ thức Chasles, đường tròn LG. Sử dụng Casio tính giá trị LG. & 1 & \textbf{Điểm, đường thẳng, mặt phẳng.} \newline \textit{Yêu cầu:} Nhận biết quan hệ liên thuộc; 3 cách xác định mặt phẳng; tìm giao tuyến, giao điểm. & 1 & \textbf{Phép biến hình và dời hình.} \newline \textit{Yêu cầu:} Nhận biết khái niệm phép dời hình, tính chất các phép đối xứng, tịnh tiến. \\ \cline{2-7}
 & 2 & (Tiếp theo) & * & & & \\ \hline

% TUẦN 2
\multirow{2}{*}{\begin{tabular}[c]{@{}c@{}}2\\(15/9-\\21/9)\end{tabular}} & 3 & \textbf{Công thức lượng giác.} \newline \textit{Yêu cầu:} Mô tả CT cộng, nhân đôi, biến đổi tích thành tổng, tổng thành tích. Giải quyết vấn đề thực tiễn. & 2 & Điểm, đường thẳng và mặt phẳng (tiếp). & 2 & Phép biến hình và dời hình (tiếp). \\ \cline{2-7}
 & 4 & (Tiếp theo) & * & & & \\ \hline

% TUẦN 3
\multirow{2}{*}{\begin{tabular}[c]{@{}c@{}}3\\(22/9-\\28/9)\end{tabular}} & 5 & \textbf{Hàm số lượng giác và đồ thị.} \newline \textit{Yêu cầu:} Nhận biết hàm chẵn, lẻ, tuần hoàn. Vẽ đồ thị $\sin, \cos, \tan, \cot$. Giải thích tập xác định, tập giá trị, chu kỳ. & 3 & Điểm, đường thẳng và mặt phẳng (tiếp). & 3 & Phép tịnh tiến. \\ \cline{2-7}
 & 6 & (Tiếp theo) & & & & \\ \hline

% TUẦN 4
\multirow{2}{*}{\begin{tabular}[c]{@{}c@{}}4\\(29/9-\\5/10)\end{tabular}} & 7 & \textbf{Phương trình lượng giác cơ bản.} \newline \textit{Yêu cầu:} Nhận biết CT nghiệm $\sin x = m, \cos x = m...$ Tính nghiệm gần đúng bằng máy tính. & 4 & Điểm, đường thẳng và mặt phẳng (tiếp). & 4 & Phép tịnh tiến. \\ \cline{2-7}
 & 8 & (Tiếp theo) & * & & & \\ \hline

% TUẦN 5
\multirow{2}{*}{\begin{tabular}[c]{@{}c@{}}5\\(6/10-\\12/10)\end{tabular}} & 9 & Bài tập cuối chương I. \textbf{\color{red}KTTX1} & 5 & \textbf{Hai đường thẳng song song.} \newline \textit{Yêu cầu:} Nhận biết vị trí tương đối (trùng, song song, cắt, chéo). & 5 & Phép đối xứng trục. \\ \cline{2-7}
 & 10 & & & & & \\ \hline

% TUẦN 9 (GIỚI HẠN)
\multirow{2}{*}{\begin{tabular}[c]{@{}c@{}}9\\(3/11-\\9/11)\end{tabular}} & 17 & \textbf{Giới hạn của dãy số.} \newline \textit{Yêu cầu:} Nhận biết khái niệm giới hạn. Giải thích giới hạn cơ bản: $\lim \frac{1}{n^k}, \lim q^n$. Tính tổng CSN lùi vô hạn. & 9 & \textbf{Hai mặt phẳng song song.} \newline \textit{Yêu cầu:} Điều kiện 2 mặt phẳng song song, định lý Thalès không gian, hình lăng trụ, hình hộp. & 9 & Phép quay. \\ \cline{2-7}
 & 18 & (Tiếp theo) & * & & & \\ \hline

% TUẦN 11 (LIÊN TỤC)
\multirow{2}{*}{\begin{tabular}[c]{@{}c@{}}11\\(17/11-\\23/11)\end{tabular}} & 21 & \textbf{Hàm số liên tục.} \textbf{\color{red}KTTX3} \newline \textit{Yêu cầu:} Nhận dạng hàm số liên tục tại điểm, trên khoảng. Xét tính liên tục của hàm sơ cấp. & 11 & \textbf{Phép chiếu song song.} \newline \textit{Yêu cầu:} Nhận biết khái niệm, tính chất phép chiếu. Vẽ hình biểu diễn khối đơn giản. & 11 & Phép vị tự. \\ \cline{2-7}
 & 22 & (Tiếp theo) & & & & \\ \hline

% THỐNG KÊ (TUẦN 17-18)
\multirow{2}{*}{\begin{tabular}[c]{@{}c@{}}17\\(30/12-\\4/1)\end{tabular}} & 29 & \textbf{Số trung bình và mốt mẫu ghép nhóm.} \newline \textit{Yêu cầu:} Tính trung bình, trung vị, tứ phân vị, mốt. Hiểu ý nghĩa đặc trưng đo xu thế trung tâm. & 31 & Trung vị và tứ phân vị của mẫu số liệu ghép nhóm. & 15 & Bài tập cuối chuyên đề 1. \\ \cline{2-3}
 & 30 & (Tiếp theo) & & & & \\ \hline

\multirow{2}{*}{\begin{tabular}[c]{@{}c@{}}18\\(5/1-\\11/1)\end{tabular}} & 32 & Trung vị và tứ phân vị (tiếp). & 34 & Thực hành Geogebra. & 16 & \textbf{\color{red}KTTX5} \newline BT cuối chuyên đề 1. \\ \cline{2-5}
 & 33 & Bài tập cuối chương V. & 35 & Dự báo dân số bằng CSN. & & \\ \hline
\end{longtable}

\vspace{1cm}
%%%%%%%%%%%
\textbf{HỌC KỲ II}
\begin{longtable}{|C{1.1cm}|C{0.5cm}|L{4.2cm}|C{0.5cm}|L{4.2cm}|C{0.5cm}|L{2.5cm}|}
\hline
\textbf{Tuần} & \textbf{T} & \textbf{ĐS-GT} & \textbf{T} & \textbf{HHKG} & \textbf{T} & \textbf{Chuyên đề} \\ \hline
\endhead

% TUẦN 1
\multirow{2}{*}{\begin{tabular}[c]{@{}c@{}}1\\(19/1-\\25/1)\end{tabular}} & 1 & \textbf{Phép tính lũy thừa.} \newline \textit{Yêu cầu:} Nhận biết lũy thừa với số mũ nguyên, hữu tỉ, thực. Giải thích các tính chất và tính giá trị biểu thức. & 1 & \textbf{Hai đường thẳng vuông góc.} \newline \textit{Yêu cầu:} Nhận biết góc giữa hai đường thẳng; chứng minh hai đường thẳng vuông góc. & 1 & \textbf{Đồ thị.} \newline \textit{Yêu cầu:} Nhận biết đường đi Euler, Hamilton. \\ \cline{2-7}
 & 2 & (Tiếp theo) & & & & \\ \hline

% TUẦN 2
\multirow{2}{*}{\begin{tabular}[c]{@{}c@{}}2\\(26/1-\\1/2)\end{tabular}} & 3 & \textbf{Phép tính logarit.} \newline \textit{Yêu cầu:} Nhận biết logarit cơ số $a$. Sử dụng tính chất logarit để rút gọn biểu thức và tính bằng máy tính. & 2 & \textbf{Đường thẳng vuông góc mặt phẳng.} \newline \textit{Yêu cầu:} Xác định điều kiện vuông góc. Tính thể tích hình chóp, lăng trụ. & 2 & Đồ thị (tiếp). \\ \cline{2-7}
 & 4 & (Tiếp theo) & & & & \\ \hline

% TUẦN 3
\multirow{2}{*}{\begin{tabular}[c]{@{}c@{}}3\\(2/2-\\8/2)\end{tabular}} & 5 & \textbf{Hàm số mũ - Hàm số logarit.} \newline \textit{Yêu cầu:} Nhận dạng và vẽ đồ thị hàm số mũ, logarit. Giải quyết bài toán thực tế. & 3 & Đường thẳng vuông góc mặt phẳng (tiếp). & 3 & Đường đi Euler và Hamilton. \\ \cline{2-7}
 & 6 & (Tiếp theo) & & & & \\ \hline

% TUẦN 4
\multirow{2}{*}{\begin{tabular}[c]{@{}c@{}}4\\(9/2-\\15/2)\end{tabular}} & 7 & \textbf{PT, bất PT mũ và logarit. \color{red}KTTX1} \newline \textit{Yêu cầu:} Giải các phương trình, bất phương trình ở dạng đơn giản. & 4 & Đường thẳng vuông góc mặt phẳng (tiếp). & 4 & Đường đi Euler và Hamilton (tiếp). \\ \cline{2-7}
 & 8 & (Tiếp theo) & & & & \\ \hline

% TUẦN 5
\multirow{2}{*}{\begin{tabular}[c]{@{}c@{}}5\\(23/2-\\1/3)\end{tabular}} & 9 & \textbf{Ôn tập chương 6.} \newline \textit{Yêu cầu:} Tổng hợp kiến thức về mũ, logarit. Giải bài toán độ pH, độ rung chấn. & 5 & \textbf{Hai mặt phẳng vuông góc.} \newline \textit{Yêu cầu:} Nhận biết 2 mặt phẳng vuông góc; tính chất lăng trụ đứng, hình hộp. & 5 & Bài toán tìm đường đi ngắn nhất. \\ \cline{2-7}
 & 10 & (Tiếp theo) & & & & \\ \hline

% TUẦN 6-7 (GIỮA KỲ)
\multirow{2}{*}{\begin{tabular}[c]{@{}c@{}}6\\(2/3-\\8/3)\end{tabular}} & 11 & \textbf{Đạo hàm. \color{red}KTTX2} \newline \textit{Yêu cầu:} Định nghĩa đạo hàm, ý nghĩa hình học (tiếp tuyến), vận dụng tính vận tốc, gia tốc. & 6 & Hai mặt phẳng vuông góc (tiếp). & 6 & Bài toán tìm đường đi ngắn nhất (tiếp). \\ \cline{2-7}
 & 12 & (Tiếp theo) & & & & \\ \hline

\textbf{7} & \multicolumn{6}{l|}{\textbf{KIỂM TRA GIỮA HỌC KỲ 2 (90 phút)}} \\ \hline

% TUẦN 8-9 (QUY TẮC ĐẠO HÀM)
\multirow{2}{*}{\begin{tabular}[c]{@{}c@{}}8\\(16/3-\\22/3)\end{tabular}} & 15 & \textbf{Các quy tắc đạo hàm.} \newline \textit{Yêu cầu:} Tính đạo hàm của tổng, hiệu, tích, thương, hàm hợp. Đạo hàm cấp hai. & 8 & \textbf{Khoảng cách trong không gian.} \newline \textit{Yêu cầu:} Khoảng cách từ điểm đến đường, mặt; khoảng cách giữa 2 đường chéo nhau. & 8 & BT cuối chuyên đề 2. \textbf{\color{red}KTTX5} \\ \cline{2-7}
 & 16 & (Tiếp theo) & & & & \\ \hline

% TUẦN 10-12
\multirow{2}{*}{\begin{tabular}[c]{@{}c@{}}10\\(31/3-\\5/4)\end{tabular}} & 19 & \textbf{Bài tập chương 7. \color{red}KTTX3} & 10 & \textbf{Góc giữa đường và mặt; Góc nhị diện.} \newline \textit{Yêu cầu:} Xác định góc, số đo góc nhị diện đơn giản. & 10 & Hình biểu diễn của một hình, khối. \\ \cline{2-7}
 & 20 & (Tiếp theo) & & & & \\ \hline

\multirow{2}{*}{\begin{tabular}[c]{@{}c@{}}11\\(6/4-\\12/4)\end{tabular}} & 21 & \textbf{Biến cố giao và quy tắc nhân xác suất.} \newline \textit{Yêu cầu:} Quy tắc cộng, nhân xác suất cho biến cố độc lập. & 11 & Ôn tập học kỳ 2. & 11 & Bản vẽ kỹ thuật. \\ \cline{2-7}
 & 22 & (Tiếp theo) & & & & \\ \hline

% CUỐI KỲ
\textbf{13-14} & \multicolumn{6}{l|}{\textbf{KIỂM TRA CUỐI HỌC KỲ 2 (TẬP TRUNG) - 90 phút}} \\ \hline

% THỰC HÀNH (TUẦN 15-16)
\multirow{2}{*}{\begin{tabular}[c]{@{}c@{}}15\\(5/5-\\10/5)\end{tabular}} & 23 & \textbf{Ứng dụng logarit đo độ pH của dung dịch.} & 13 & Vẽ hình khối bằng phần mềm Geogebra. & 13 & BT cuối chuyên đề 3. \\ \cline{2-5}
 & & & 14 & Làm kính 3D để quan sát. & & \\ \hline

\multirow{2}{*}{\begin{tabular}[c]{@{}c@{}}16\\(11/5-\\17/5)\end{tabular}} & 24 & (Tiếp theo thực hành pH) & 15 & Vẽ hình khối bằng Geogebra (tiếp). & 14 & BT cuối chuyên đề 3 (tiếp). \\ \cline{2-5}
 & & & 16 & & & \\ \hline

% LUYỆN TẬP CUỐI NĂM
\textbf{17} & * & \textbf{Luyện tập tổng hợp} & * & \textbf{Luyện tập tổng hợp} & * & \textbf{Luyện tập tổng hợp} \\ \hline
\end{longtable}

\vspace{1cm}
% --- CHỮ KÝ ---
\noindent
\begin{minipage}[t]{0.45\textwidth}
    \centering
    \textbf{P. HIỆU TRƯỞNG}\\
    \vspace{2.5cm}
    \textbf{...........................................}
\end{minipage}
\hfill
\begin{minipage}[t]{0.45\textwidth}
    \centering
    \textit{Tp. HCM, ngày 05 tháng 09 năm 2025}\\
    \textbf{TỔ TRƯỞNG}\\
    \vspace{2cm}
    \textbf{Đỗ Thị Bạch Lan}
\end{minipage}

\end{document}